Una masa de 100 kg cuelga de un resorte de 80 N/m y un amortiguador viscoso de 40 N·s/m, dispuestos en paralelo y fijados a un soporte inmóvil. Mide x hacia abajo desde el equilibrio previo bajo el peso. Una fuerza adicional de 15 N actúa hacia abajo durante 8 s y luego se retira. En \(t=0\), x=x\ensuremath{^{\prime}}=0. Se suponen componentes lineales y parámetros constantes.

\begin{enumerate}[label=(\alph*),ref=\alph*,leftmargin=*]
\item Formula la ecuación diferencial.
\item Obtén x(t).
\item Dibuja la gráfica.
\end{enumerate}
